\documentclass[11pt]{article}

% Load your custom style file.
% Adjust the path if needed depending on where this .tex file lives.
\usepackage[solutions]{Style}

\begin{document}

% Set the homework information that appears at the top of the page.
% Order: student name, course, assignment title, due date, collaborators.
\sethwinfo{Hem Dadlani}{Math 55}{Sample Homework Write-Up}{April 13, 2026}[None]

% Print the formatted homework title block.
\makehwtitle

% A section header. Problem numbers will reset by section because of Style.sty.
\section{Linear Algebra Warm-Up}

% Example problem environment.
% First optional argument = problem title.
% Second optional argument = points.
\begin{problem}[A basic proof][10]
    Prove that if $v \in \R^n$ and $0v = 0$, then $v + 0 = v$.
\end{problem}

% The solution environment only appears because we loaded Style.sty with [solutions].
% If you switch to [nosolutions], everything inside solution environments is hidden.
\begin{solution}
    By the vector space axioms, $0v = 0$ for every vector $v$, and $0$ is the additive identity.
    Hence
    \[
        v + 0 = v.
    \]
\end{solution}

% Example with subproblems.
\begin{problem}[Quick computations][12]
    Let
    \[
        A = \begin{pmatrix}1 & 2\\ 3 & 4\end{pmatrix}.
    \]
    Compute the following.
    \begin{subproblem}
        \item $\Tr(A)$
        \item $\det(A)$
        \item $A\transpose$
    \end{subproblem}
\end{problem}

\begin{solution}
    \begin{subproblem}
        \item Since $\Tr(A)$ is the sum of the diagonal entries,
        \[
            \Tr(A) = 1+4=5.
        \]
        \item We compute
        \[
            \det(A)=1\cdot 4 - 2\cdot 3 = -2.
        \]
        \item The transpose is
        \[
            A\transpose = \begin{pmatrix}1 & 3\\ 2 & 4\end{pmatrix}.
        \]
    \end{subproblem}
\end{solution}

% Example theorem-style environments from Style.sty.
\begin{definition}[Span]
    If $v_1,\dots,v_k \in V$, then
    \[
        \Span(v_1,\dots,v_k)
    \]
    is the set of all linear combinations of $v_1,\dots,v_k$.
\end{definition}

\begin{remark}
    The command \verb|\Span| is already defined in your style file as a math operator,
    so it gets the correct spacing automatically.
\end{remark}

\begin{theorem}[Tiny example]
    Let $V$ be a vector space and let $v \in V$. Then $\Span(v)$ is a subspace of $V$.
\end{theorem}

\begin{proof}
    To show $\Span(v)$ is a subspace, check the usual three properties:
    it contains $0$, it is closed under addition, and it is closed under scalar multiplication.
    Indeed, every element of $\Span(v)$ has the form $av$ for some scalar $a$.
    If $av,bv \in \Span(v)$, then
    \[
        av+bv=(a+b)v \in \Span(v),
    \]
    and for any scalar $c$,
    \[
        c(av)=(ca)v \in \Span(v).
    \]
    Also, $0v=0 \in \Span(v)$.
\end{proof}

\section{Using the custom commands}

% This problem shows off several macros from Style.sty.
\begin{problem}[Style.sty macro demo][15]
    Let $x,y \in \R^n$. Write a short display that uses the custom commands
    \verb|\abs|, \verb|\norm|, \verb|\inner|, \verb|\set|, and \verb|\ceil|.
\end{problem}

\begin{solution}
    Here are a few examples:
    \[
        \abs{-7} = 7,
        \qquad
        \norm{x+y} \le \norm{x} + \norm{y},
        \qquad
        \inner{x}{y} = \sum_{i=1}^n x_i y_i.
    \]
    \[
        S = \set*{n \in \N \given n \text{ is even}},
        \qquad
        \ceil*{\frac{7}{2}} = 4.
    \]

    % The starred versions like \ceil* or \norm* automatically size the delimiters.
    The starred delimiter forms automatically adjust the size of the brackets when needed.
\end{solution}

% Utility command examples.
\paragraph{A few utility commands}

This is a normal sentence. \highlight{This text is highlighted using the custom command.}

\TODO{Replace this sample TODO with something real before submitting.}

% marginref puts a note in the margin.
\marginref{Useful when reviewing drafts.}

\end{document}